%StructureSemantics.tex

\ksec{Structures Semantics (KerML 8.4.4.4)}

\pn \texttt{Structure} is \texttt{Class}, but structures implicitly specialize \texttt{Objects::Object}.  \texttt{Object} is disjoint from \texttt{Performance}.
% yet has three features (or steps) which are zero or more  \texttt{Performance}.

\pn The example in KerML \S8.4.4.4 has recursion which seems wrong.  Is this expressing an infinite binary tree of \texttt{S}?%\footnote{\texttt{Object::objects} s.b. \texttt{Objects::objects}}
\ecaption{Recursive Structure}
\begin{lstlisting}
struct S specializes Objects::Object {
	feature a : S subsets Objecst::objects;
	composite feature b : S subsets Objects::Object::subobjects;
}
\end{lstlisting}

\pn An \texttt{Object::Object} may have \texttt{subobjects} which are defined using metaphysical mereology (\ref{eq:df-par}).
\ecaption{subobjects}
\begin{lstlisting}
	abstract struct Object specializes Occurrence {
		doc /* Object is the most general class of structural occurrences that may change over time. */
		feature self: Object redefines Occurrence::self;
		
		composite feature subobjects: Object[0..*] subsets suboccurrences {
			doc /* The suboccurrences of this Object that are also Objects. */
			language "Domain" /* forall so in subobjects are Mereology::PartOf(so,self) */
		}
\end{lstlisting}

\pn All subobjects of an objects are parts of that object (\ref{eq:df-par}).\footnote{
\texttt{subsets} \texttt{objects} removed because it was redundant.}

%\begin{restatable}{definition}{subobjects}~\blTitle{df-subobjects}~ Define \texttt{subobjects}.  
%\begin{align}
%\forall s \in \texttt{subobjects}~\vert~s~P~\texttt{self}\label{eq:df-subobjects}\tag{df-subobjects}
%\end{align}
%\end{restatable}


\pn An \texttt{Object} may be either physical or virtual depending on whether it has a \texttt{structuredSpaceBoun\-dary} which is its
\texttt{RegionSurface} of its \texttt{Location} (\ref{eq:df-lfu}).\footnote{
Reconciliation of \texttt{Regions::Surface} with \texttt{structuredSpaceBoundary} yet to be done.  Another case where \texttt{portion} may be incorrect.}

\ecaption{structuredSpaceBoundary}
\begin{lstlisting}
		portion structuredSpaceBoundary : StructuredSpaceObject[0..1] subsets spaceBoundary {
			doc
			/*
			 * A space boundary that is a structured space object.
			 */
			language "Domain" /* not structuredSpaceBoundary->empty implies
			  * Regions::RegionSurface(Regions::Location(self)) = structuredSpaceBoundary */
		}
\end{lstlisting}

%\begin{restatable}{definition}{location}~\blTitle{df-location}~ Define object location.  
%\begin{align}
%\texttt{structuredSpaceBoundary} \ne \emptyset~\rightarrow~\texttt{self}~L~\texttt{structuredSpaceBoundary}\label{eq:df-location}\tag{df-location}
%\end{align}
%\end{restatable}

\pn Physical objects may only contain other physical objects.  Virtual objects are allocated to physical objects.  By expansivity (\ref{eq:df-pex}), parthood implies regional containment.

%\begin{lstlisting}
%		composite feature subobjects: Object[0..*] subsets objects, suboccurrences
%			intersects objects, suboccurrences {
%			doc
%			/*
%			 * The suboccurrences of this Object that are also Objects.
%			 */
%		}
%\end{lstlisting}

\begin{restatable}{theorem}{containment}~\blTitle{bl.containment}~ An object's region contains the regions of its subobjects.
\begin{align}
\vdash \texttt{structuredSpaceBoundary} \ne \emptyset~\rightarrow~%\notag \\
\forall s \in \texttt{subobjects}~\vert~\notag \\
s.\texttt{structuredSpaceBoundary} \ne \emptyclass \rightarrow \notag \\
\texttt{structuredSpaceBoundary}~C~s.\texttt{structuredSpaceBoundary}\label{eq:bl.containment}\tag{bl.containment}
\end{align}
\end{restatable}

\pn Structures are not behaviors, but they may have behaviors.  Any behavior must specialize \texttt{Performances::}\-\texttt{Performance}. 
Structures may have three kinds of \texttt{Performance},  involving, enacted, and owned.

\pn Involving just says this object is ``involved'' in the performance, perhaps in concert with other objects.
\ecaption{involvingPerformances}
\begin{lstlisting}
		feature involvingPerformances: Performance[0..*] subsets performances {
			doc
			/*
			 * Performances in which this object is involved.
			 */
		}
\end{lstlisting}

\pn Enacted performances during the existence of the object.
\ecaption{enactedPerformances}
\begin{lstlisting}
		abstract step enactedPerformances: Performance[0..*] subsets involvingPerformances {
			doc
			/*
			 * Performances that are enacted by this object.
			 */
			language "Domain" /* forall ep in enactedPerformances are self->timeEnclosedOccurrence(ep) */
		}
\end{lstlisting}

\pn Owned performances are contained and performed by the object.
\ecaption{ownedPerformances}\label{ownedPerformances}
\begin{lstlisting}
		composite step ownedPerformances: Performance[0..*] 
		  subsets involvingPerformances, suboccurrences
			intersects involvingPerformances, suboccurrences {
			language "Domain" /* forall op in ownedPerformances are self->
			  * (timeEnclosedOccurrence(ep) and Mereology::PartOf(op,self) */
			doc
			/*
			 * Performances that are owned by this object.
			 */
			 
			feature redefines this default that {
				doc
				/*
				 * The owning object is the default "this" reference for all ownedPerformances.
				 */
			}
		}
\end{lstlisting}

\pn \texttt{Objects} defines link-objects which are both associations and structures.

\pn \texttt{Body}, \texttt{Surface}, and \texttt{Point} have close correlation with \texttt{Regions}' \texttt{Region}, \texttt{Surface} and \texttt{Point}.
There is no  \texttt{Regions} equivalent for  \texttt{Curve}.

%\pn  \texttt{StructuredSpaceObject} is a mess because it tries


